\section{Established conclusions, excluded candidates, and open questions}
\label{sec:status-literature}

\paragraph{Status of this research branch.}
All statements in this appendix are presented as working proofs and
computational certificates under verification. Internal checks of identities
and repeated applications of published criteria are not external independent
peer review. Below, the word \emph{consequence} denotes a logical deduction
from the displayed formulas and cited theorems; it does not assert priority
in the literature.

\begin{center}
\small
\begin{tabular}{@{}p{0.30\linewidth}p{0.64\linewidth}@{}}
\toprule
Question & Conclusion and its precise scope\\
\midrule
The specific Shastri curve &
A polynomial rectifier over $K=\Q(\sqrt{-3})$ excludes this
parametrization as a candidate for complex nonrectifiability. It does
not exclude other nonrectifiable spatial lines.\\[4pt]
The harmonic model and $C_N$ &
The specified parametrizations are rectified. The harmonic model
uses $D_5=t^5-5t^3+5t$; replacing $5t$ by $4t$ is not covered
by this statement.\\[4pt]
Tameness of the constructed maps &
Wildness of $\Phi_T,\Phi_H,\Phi_N$ follows from relative-coordinate
criteria. The nonexistence of \emph{any} tame rectifier of the same
curves has not been proved.\\[4pt]
Stabilization &
$\Phi_\ast\times\id_{\A^3}\in\Tame_6(K)$ follows from a known
stable-tameness theorem. The minimum number of additional variables
has not been determined.\\[4pt]
Cremona subgroup &
The constructed maps belong to $\mathcal C_3(K)$.
This neither describes the whole group $\Cr_3(K)$ nor provides the
first example of a wild element of $\mathcal C_3(K)$.\\[4pt]
Field of definition &
For Shastri's curve, the minimum degree of a suitable number field
over $\Q$ is $2$. Uniqueness of a suitable quadratic field and an
effective classification of rectification fields are not established.\\
\bottomrule
\end{tabular}
\end{center}

\paragraph{The historical candidate and the general problem.}
Shastri's paper \cite[\S1 and \S3, Example~1]{Shastri1992} proposes
real knot parametrizations as candidates for investigating complex
nonrectifiability. Thus a complete polynomial automorphism for
$\gamma_T$ answers the rectification question for \emph{this curve}
affirmatively. It does not disprove the historical conjecture that at
least one nonrectifiable embedding
$\A^1_{\C}\hookrightarrow\A^3_{\C}$ exists. Nor does it prove the
opposite universal conjecture that every such embedding is rectifiable.
We do not attribute to Shastri's paper a separate numbered conjecture
that it does not state.

\paragraph{Related conjectures from the first appendix.}
The consequences for the van den Essen--van Rossum candidates already
recorded in the first appendix must be accounted for separately
\cite[\S5, Conjectures~5.1--5.3]{VdER2002}.\footnote{The numbering and
identification of the objects here are inherited from Appendix 1.
The publisher's record was checked in preparing this version, but a
fresh download of the full AMS text was unavailable. The algebraic
deductions below were checked directly; this limitation of the
bibliographic audit is also recorded in
\texttt{audit/claim\_status.md}.}
These are consequences of the original rectification, rather than
new results about the Cremona group. For $\gamma_T=(f,g,h)$, consider
\[
 D=t\partial_u+f'(u)\partial_x+g'(u)\partial_y+h'(u)\partial_z
 \quad\text{on }\C[t,u,x,y,z].
\]
Conjecture~5.1, also called \emph{Conjecture E} in subsequent
literature, predicted that $\ker D$ is not a polynomial algebra.
From the rectifier $\Phi_T=(P,Q,R)$ with inverse $\Psi$, setting
$V=\gamma_T(u)-t(x,y,z)$ gives
\[
 p=P(V),\qquad q=Q(V)/t,\qquad r=R(V)/t,\qquad
 s=(u-P(V))/t.
\]
All divisions are exact because $\Phi_T\gamma_T=(u,0,0)$.
The inverse change of variables is
\[
 u=p+ts,\qquad
 (x,y,z)=\frac{\gamma_T(p+ts)-\Psi(p,tq,tr)}{t};
\]
it is polynomial as well. Consequently,
\[
 \C[t,u,x,y,z]=\C[t,p,q,r,s],\qquad
 D=\partial_s,\qquad \ker D=\C[t,p,q,r].
\]
This negates precisely the specified conjecture~5.1. If $\sigma$
denotes the original slice in the same construction, then
$\sigma-s\in\ker D$, whence
$\C[t,u,x,y,z]/(\sigma)\simeq\C^{[4]}$.
Furthermore, for $I_{\gamma_T}=(Q,R)$ we have
\[
 \C[x,y,z][t,I_{\gamma_T}/t]
 =\C[t,P,Q/t,R/t]\simeq\C^{[4]}.
\]
These are the corresponding polynomial descriptions of the objects
in~5.2 and~5.3. Detailed formulas, identification of the notation,
and certificates are provided in the repository's first appendix.
These particular calculations do not solve the general cancellation,
embedding, or action-linearization problems. Oda's withdrawn preprint
\cite{OdaWithdrawn} is not used as a proof: its arXiv record states
that it was withdrawn on 9 May 2020 because of an error.

\paragraph{Predecessors of the programme.}
The connection between the embedding problem and Cremona equivalence
is explicitly discussed by Mella--Polastri
\cite[introduction and Theorem~1]{MP2009}.
Relative coordinates, factorizations over the fraction field, and
transfer are studied by Edo--V\'en\'ereau \cite{EV2001}.
The closest predecessor of the constructive part is Lewis
\cite[Lemma~1, Theorems~2 and~12]{Lewis2024}: a coordinate restricting
to the parameter is constructed using residual coordinates and
corrections of rational transformations. That paper also compares
earlier results of Craighero, Bhatwadekar--Roy, and Kuroda.
Accordingly, the general procedure \emph{localize, then eliminate
denominators} is not claimed here as a new strategy.

The existence of wild automorphisms in three variables has long been
established; the passage from a wild relative coordinate to wildness
in three variables is supplied by \cite{UY2004}. Blanc--Hed\'en
already proved $\operatorname{Nagata}\in\mathcal C_3$
\cite[Proposition~6.8]{BH2015}. Stable tameness for the class used
here follows from \cite[Corollary~4.9]{BEW2012}.
The contribution of this appendix consists in applying these results
to the \emph{specific} rectifiers, verifying their hypotheses, and
formulating an explicit programme for further work.

\paragraph{Unresolved questions and limits of the search.}
Questions remaining open within this project include completeness
of the relative method, the existence of a tame rectifier for the
same curves, minimum rectifier degrees, descent to other totally
imaginary number fields, and the general case
$\A^1\hookrightarrow\A^3$ over $\C$.
Finite computations for $N\leq N_{\max}$ check the implementation;
statements for every $N$ rely on inductive proofs.
Failure of a bounded search does not prove nonrectifiability.
A targeted literature search as of 28 October 2026 did not identify
an earlier publication containing our complete Shastri rectifier;
failure to find such a publication does not establish priority.
